\paragraph{分量余类的压力函数与素数方向饱和}
沿用稳定类型 \(x\in X_m\) 的候选位点诱导图分解
\[
\Gamma_x \cong \bigsqcup_{j=1}^{r(x)} P_{\ell_j(x)},
\]
以及纤维多重度的 Fibonacci 乘积分解
\[
d_m(x):=\abs{\Fold_m^{-1}(x)}=\prod_{j=1}^{r(x)} F_{\ell_j(x)+2}.
\]
内部文档 \cite{Internal2026FoldComponentResiduePressurePAdicSaturation20260302} 在模 \(3\) 的分量计数中给出有限状态在线归一化；下述结果给出一般模数的压力函数表达，并由此导出素数方向的指数型饱和与典型光滑性。

\begin{theorem}[任意模数下的分量余类计数：压力函数与指数型下尾]\label{thm:fold-component-residue-pressure-ld}
取整数 \(L\ge 2\) 与同余类 \(a\in \ZZ/L\ZZ\)。对 \(x\in X_m\) 定义分量余类计数
\[
N_{m}^{(L,a)}(x):=\#\{j:\ \ell_j(x)\equiv a\ (\mathrm{mod}\ L)\}.
\]
在 \(X_m\) 上取均匀分布 \(x\sim\mathrm{Unif}(X_m)\)。则存在常数 \(\mu_{L,a}>0\)、\(c_{L,a}>0\)，使得对任意固定 \(\varepsilon\in(0,\mu_{L,a})\)，当 \(m\) 充分大时
\[
\PP_{x\sim\mathrm{Unif}(X_m)}\!\left[N_{m}^{(L,a)}(x)\le \varepsilon m\right]\le e^{-c_{L,a}m}.
\]
并且对数矩母函数极限存在且可由 Perron 根表示：存在有限状态有向图 \(G_{L}\) 与边可观测 \(\varphi_{L,a}:E(G_L)\to\{0,1\}\)，使得对任意 \(x\in X_m\)，若以 \(\gamma_1(x),\dots,\gamma_m(x)\in E(G_L)\) 表示其对应的长度 \(m\) 状态路径，则
\[
N_{m}^{(L,a)}(x)=\sum_{t=1}^{m}\varphi_{L,a}(\gamma_t(x))+O(1),
\]
并令 \(B_{\theta}\) 为按权 \(e^{\theta\varphi_{L,a}}\) 加权的邻接矩阵，则
\[
\Lambda_{L,a}(\theta)
:=\lim_{m\to\infty}\frac1m\log \EE_{x\sim\mathrm{Unif}(X_m)}\!\left[e^{\theta N_m^{(L,a)}(x)}\right]
=\log \rho(B_{\theta})-\log \rho(B_{0}),
\]
其中 \(\rho(\cdot)\) 为谱半径。特别地，\(\Lambda_{L,a}\) 在 \(\RR\) 上实解析，且 \(\mu_{L,a}=\Lambda_{L,a}'(0)>0\)。
\end{theorem}
\begin{proof}
内部文档 \cite{Internal2026FoldComponentResiduePressurePAdicSaturation20260302} 在 \(L=3\) 的情形中构造了“有限状态在线归一化”：逐位读入 \(x\in X_m\) 的同时，以有限状态扩张同步跟踪 \(\Gamma_x\) 的分量边界，并记录当前分量长度的模 \(3\) 信息，从而把余类计数写成有限状态路径上的加法泛函，并得到相应的加权邻接矩阵与压力函数表示。

将该构造中的模 \(3\) 计数器替换为模 \(L\) 计数器（即把状态空间并入 \(\ell\mapsto \ell+1\ (\mathrm{mod}\ L)\) 的有限记忆），得到所述 \(G_L\) 与 \(\varphi_{L,a}\)，并使
\(
N_{m}^{(L,a)}(x)=\sum_{t=1}^{m}\varphi_{L,a}(\gamma_t(x))+O(1)
\)
成立。由正权矩阵的 Perron--Frobenius 理论，\(\rho(B_\theta)\) 严格为正且对 \(\theta\) 实解析，故 \(\Lambda_{L,a}\) 的极限表达成立并实解析。
由于 \(\varphi_{L,a}\) 非恒零且状态核强连通，得到 \(\mu_{L,a}=\Lambda_{L,a}'(0)>0\)。最后由 Chernoff 界（等价于 \(\Lambda_{L,a}\) 的凸性与 Legendre 对偶）得到任意 \(\varepsilon\in(0,\mu_{L,a})\) 的指数型下尾界。证毕。
\end{proof}
\leanverified{paper\_fold\_component\_residue\_pressure\_ld}

\leanverified{paper\_fold\_padic\_saturation\_prime}
\begin{corollary}[素数方向的线性 \(p\)-进饱和]\label{cor:fold-padic-saturation-prime}
对素数 \(p\) 定义 Fibonacci 数列模 \(p\) 的出现阶
\[
Z(p):=\min\{n\ge 1:\ p\mid F_n\}.
\]
则对任意 \(n\ge 1\) 有整除判别
\[
p\mid F_n\ \Longleftrightarrow\ Z(p)\mid n.
\]
令
\[
a_p\equiv -2\pmod{Z(p)},\qquad N_{m,p}(x):=N_m^{(Z(p),a_p)}(x).
\]
则存在常数 \(\alpha_p>0\)、\(c_p>0\)，使得当 \(m\) 充分大时
\[
\PP_{x\sim \mathrm{Unif}(X_m)}\!\left[v_p(d_m(x))\le \alpha_p m\right]\le e^{-c_p m},
\]
其中 \(v_p\) 为 \(p\)-进赋值。特别地，
\[
\PP_{x\sim \mathrm{Unif}(X_m)}[\,p\mid d_m(x)\,]\ \ge\ 1-e^{-c_p m}\xrightarrow[m\to\infty]{}1.
\]
\end{corollary}
\begin{proof}
由 \(\gcd(F_a,F_b)=F_{\gcd(a,b)}\) 可知集合 \(\{n:\ p\mid F_n\}\) 在最大公因子下闭合，从而由出现阶定义立即推出 \(p\mid F_n\Rightarrow Z(p)\mid n\)；反向蕴含由定义显然成立。

若 \(\ell_j(x)\equiv a_p\ (\mathrm{mod}\ Z(p))\)，则 \(\ell_j(x)+2\equiv 0\ (\mathrm{mod}\ Z(p))\)，由整除判别得 \(p\mid F_{\ell_j(x)+2}\)。由乘积分解
\(
d_m(x)=\prod_{j=1}^{r(x)}F_{\ell_j(x)+2}
\)
得到
\[
v_p(d_m(x))=\sum_{j=1}^{r(x)} v_p\!\bigl(F_{\ell_j(x)+2}\bigr)\ \ge\ \#\{j:\ p\mid F_{\ell_j(x)+2}\}\ \ge\ N_{m,p}(x).
\]
对 \(N_{m,p}(x)\) 应用定理 \ref{thm:fold-component-residue-pressure-ld}（取 \(L=Z(p)\)、\(a=a_p\)），存在 \(\mu_p:=\mu_{Z(p),a_p}>0\)，并可取任意 \(\alpha_p\in(0,\mu_p)\) 使得当 \(m\) 充分大时
\[
\PP\!\left[N_{m,p}(x)\le \alpha_p m\right]\le e^{-c_p m}.
\]
结合 \(v_p(d_m(x))\ge N_{m,p}(x)\) 得第一式；并由事件包含
\(
\{p\nmid d_m(x)\}\subseteq \{v_p(d_m(x))\le \alpha_p m\}
\)
得到整除概率的指数收敛。证毕。
\end{proof}

\begin{remark}\label{rem:fold-2adic-saturation-mod3}
当 \(p=2\) 时 \(Z(2)=3\) 且 \(a_2\equiv 1\pmod 3\)。因此推论 \ref{cor:fold-padic-saturation-prime} 的 \(p=2\) 特化等价于模 \(3\) 余类 \(1\) 的线性出现，并与定理 \ref{thm:killo-fold-fiber-homotopy-parity-equivalence} 的奇偶判别与同伦二分相容。
\end{remark}

\begin{corollary}[有限素数集的同步线性饱和]\label{cor:fold-padic-saturation-finite-set}
\leanverified{paper\_fold\_padic\_saturation\_finite\_set}
取有限素数集 \(\mathcal P\subset\mathbb P\)。则存在常数 \(\{\alpha_p>0\}_{p\in\mathcal P}\) 与 \(c>0\)，使得当 \(m\) 充分大时
\[
\PP_{x\sim \mathrm{Unif}(X_m)}\!\Big[\ \forall p\in\mathcal P,\ v_p(d_m(x))\ge \alpha_p m\ \Big]\ \ge\ 1-e^{-c m}.
\]
\end{corollary}
\begin{proof}
对每个 \(p\in\mathcal P\) 应用推论 \ref{cor:fold-padic-saturation-prime}，并取 \(\alpha_p\) 与 \(c_p\) 相应常数，则
\(
\PP[v_p(d_m(x))<\alpha_p m]\le e^{-c_p m}
\)。
由并集界得
\[
\PP\bigl[\exists p\in\mathcal P:\ v_p(d_m(x))<\alpha_p m\bigr]
\le \sum_{p\in\mathcal P} e^{-c_p m}
\le e^{-c m}
\]
（取 \(c=\frac12\min_{p\in\mathcal P}c_p\) 并令 \(m\) 足够大）。证毕。
\end{proof}

\begin{theorem}[典型纤维的多项式光滑性]\label{thm:fold-typical-fiber-polynomial-smoothness}
\leanverified{paper\_fold\_typical\_fiber\_polynomial\_smoothness}
记
\[
\ell_{\max}(x):=\max_j \ell_j(x)\quad(\text{若无分量则取 }0),
\qquad
P^+(n):=\max\{p\in\mathbb P:\ p\mid n\},
\]
并约定 \(P^+(1)=1\)。则存在常数 \(A>0\) 与 \(C,c>0\)，使得对所有充分大 \(m\) 有
\[
\PP_{x\sim \mathrm{Unif}(X_m)}\!\left[\ell_{\max}(x)\ge A\log m\right]\ \le\ C\,m\,e^{-cA\log m}=O(m^{1-cA}),
\]
从而选取 \(A>1/c\) 得
\[
\ell_{\max}(x)=O(\log m)\quad\text{以概率 }1-o(1).
\]
进一步，存在常数 \(\kappa>0\) 使得
\[
\PP_{x\sim \mathrm{Unif}(X_m)}\!\left[P^+(d_m(x))\le m^{\kappa}\right]\xrightarrow[m\to\infty]{}1.
\]
\end{theorem}
\begin{proof}
对稳定类型的候选位点诱导图分解，有限状态在线归一化不仅可跟踪模 \(L\) 信息，也可只跟踪“当前是否仍处于同一分量内部”的内部状态集合。令 \(G\) 为对应有限有向图，并令 \(S\subset V(G)\) 为“处于某一分量内部且尚未结束”的状态集合。由分量在有限步内可结束，\(S\) 为真子系统，故其诱导子图满足
\(
\rho(G|_S)<\rho(G)
\)。
因此存在常数 \(\eta\in(0,1)\)、\(K>0\)，使得任意给定起点下，单个分量长度满足几何尾界
\[
\PP[\text{某分量长度}\ge t]\le K\,\eta^{t}.
\]
一条长度 \(m\) 的路径最多出现 \(m\) 个分量，故并集界给出
\[
\PP[\ell_{\max}(x)\ge t]\le Km\,\eta^{t}.
\]
取 \(t=A\log m\) 即得第一段结论。

由 Fibonacci 乘积表达式，
\[
P^+(d_m(x))\le d_m(x)\le \prod_{j=1}^{r(x)}F_{\ell_j(x)+2}\le \Bigl(\max_j F_{\ell_j(x)+2}\Bigr)^{r(x)}.
\]
更直接地，
\[
P^+(d_m(x))\le \max_j P^+\!\bigl(F_{\ell_j(x)+2}\bigr)\le \max_j F_{\ell_j(x)+2}=F_{\ell_{\max}(x)+2}.
\]
由 \(F_n\le C_0\varphi^{n}\)（\(\varphi\) 为黄金比）与 \(\ell_{\max}(x)=O(\log m)\) 以概率 \(1-o(1)\) 成立，存在 \(\kappa>0\) 使得在该高概率事件上
\[
F_{\ell_{\max}(x)+2}\le m^\kappa.
\]
从而 \(\PP[P^+(d_m(x))\le m^\kappa]\to 1\)。证毕。
\end{proof}

\begin{corollary}[“饱和—光滑”二象性]\label{cor:fold-saturation-smooth-dichotomy}
对任意有限素数集 \(\mathcal P\subset\mathbb P\)，存在 \(\kappa>0\)、\(\{\alpha_p\}_{p\in\mathcal P}\) 与 \(c>0\)，使得当 \(m\) 充分大时
\[
\PP_{x\sim\mathrm{Unif}(X_m)}\!\Big[\ P^+(d_m(x))\le m^{\kappa}\ \text{且}\ \forall p\in\mathcal P,\ p^{\alpha_p m}\mid d_m(x)\ \Big]\ \ge\ 1-e^{-c m}.
\]
\end{corollary}
\begin{proof}
由定理 \ref{thm:fold-typical-fiber-polynomial-smoothness} 控制最大素因子，由推论 \ref{cor:fold-padic-saturation-finite-set} 控制有限素数集的线性赋值；取二者事件交集并以并集界估计失败概率即得。证毕。
\end{proof}

\begin{conjecture}[分量余类的渐近等分布与出现阶普适律]\label{conj:fold-component-residue-equidistribution}
设分量密度极限
\(
\nu:=\lim_{m\to\infty} r(x)/m
\)
在 \(x\sim\mathrm{Unif}(X_m)\) 下以概率收敛的意义存在。则对每个固定 \(L\ge 2\)，余类计数向量满足
\[
\frac1m\bigl(N_m^{(L,a)}(x)\bigr)_{a\in\ZZ/L\ZZ}\ \xrightarrow[m\to\infty]{\PP}\ \Bigl(\frac{\nu}{L},\dots,\frac{\nu}{L}\Bigr).
\]
由此将推出推论 \ref{cor:fold-padic-saturation-prime} 中的线性系数满足
\[
\alpha_p \asymp \frac{\nu}{Z(p)}\qquad(p\ \text{素}),
\]
从而出现阶 \(Z(p)\) 在一阶上控制 \(p\)-进饱和强度。
\end{conjecture}
